\section{Young tableaux, hook products and insertion}
\label{sec:tableaux}

The overview article links to \emph{Young diagram}, which redirects to
\emph{Young tableau}, and illustrates the partition \((5,4,1)\). The
linked article gives both hook product formulas and numerical examples
for this shape. We derive the counts directly, and use insertion to
connect them with permutations.

A shape is a list \(\lambda_1\ge\cdots\ge\lambda_r\ge0\) of sum \(n\),
with trailing zeros allowed for a fixed row bound. In a sum over partitions
we omit trailing zeros, so each shape occurs once. Its boxes are \((i,j)\), \(1\le i\le r\),
\(1\le j\le\lambda_i\). A standard Young tableau fills the boxes
bijectively with \(1,\ldots,n\), strictly increasing along rows and
down columns. A semistandard tableau on \(1,\ldots,r\) is weakly
increasing along rows and strictly increasing down columns. The hook
length \(h_{ij}\) is one plus the number of boxes strictly to the
right in its row or strictly below in its column.

\begin{center}
\begin{minipage}{.32\linewidth}\centering
\begin{tikzpicture}[x=.41cm,y=-.41cm,baseline]
\foreach \i/\j/\h in {0/0/7,1/0/5,2/0/4,3/0/3,4/0/1,0/1/5,1/1/3,2/1/2,3/1/1,0/2/1}{\draw[draw=inkblue,fill=soft](\i,\j)rectangle ++(1,1);\node[font=\footnotesize\sffamily] at (\i+.5,\j+.5) {\h};}
\end{tikzpicture}
\end{minipage}\begin{minipage}{.58\linewidth}\small
Each box shows its hook length. The product for shape $(5,4,1)$ is $12{,}600$, so the theorem below gives $10!/12{,}600=288$ standard fillings.
\end{minipage}
\end{center}

\subsection{Corner recursion and the hook-length formula}

\begin{theorem}[Frame--Robinson--Thrall hook-length formula]\label{result15}
If \(f^\lambda\) counts the standard tableaux of shape \(\lambda\), then
\[
 f^\lambda \prod_{(i,j)\in\lambda} h_{ij}=n!.
\]
Transposition preserves \(f^\lambda\). For a rectangle with \(r\) rows
of length \(s\), the count is \((rs)!\prod_{j=0}^{r-1} j!/(s+j)!\).
For two rows of lengths \(a\ge b\) it is
\(\binom{a+b}b-\binom{a+b}{b-1}\), interpreting the second term as zero
for \(b=0\).
\end{theorem}
\begin{proof}
Put \(x_i=\lambda_i+r-i\), so \(x_1>\cdots>x_r\ge0\), and write
\(\Delta(x)=\prod_{i<j}(x_i-x_j)\). The hooks in row \(i\) are exactly
the numbers \(1,\ldots,x_i\) with the \(r-i\) numbers \(x_i-x_k\)
for \(k>i\) omitted. Indeed the row's hook lengths
\[
 h_{ij}=\lambda_i-j+1+\#\{k>i:\lambda_k\ge j\}
\]
strictly decrease with \(j\) and lie in that range. For \(k>i\), if
\(j\le\lambda_k\) the lower rows through \(k\) all contribute, giving
\(h_{ij}>x_i-x_k\). If \(j>\lambda_k\), at most \(k-i-1\) lower rows
contribute and \(h_{ij}<x_i-x_k\). The omitted numbers are distinct,
and \(x_i-(r-i)=\lambda_i\) accounts for all others. Multiplying the row
products gives
\[
 \prod_{(i,j)\in\lambda}h_{ij}=\frac{\prod_i x_i!}{\Delta(x)}.
\]

For \(n>0\), the largest tableau entry occupies a removable corner,
namely the last box in some row \(i\) with
\(\lambda_i>\lambda_{i+1}\), where \(\lambda_{r+1}=0\).
Deleting it and reattaching \(n\) give inverse bijections. Thus
\(f^\lambda=\sum_{i\ {\rm removable}} f^{\lambda-e_i}\) with
\(f^\varnothing=1\). We verify this recursion for the positive rational
number \(W(\lambda)=n!\Delta(x)/\prod_i x_i!\). For a removable row,
\[
 \frac{nW(\lambda-e_i)}{W(\lambda)}=x_i A_i,\qquad
 A_i=\prod_{j\ne i}\frac{x_i-x_j-1}{x_i-x_j}.
\]
This follows by dividing the two factorial products and the factors of
\(\Delta\) involving \(i\). A nonremovable row contributes \(x_iA_i=0\):
either \(i<r\) and \(x_i-x_{i+1}=1\), or \(i=r\) and \(x_i=0\).

We need \(\sum_i x_iA_i=\sum_i x_i-\binom r2=n\). Here is a finite
algebraic proof. Put \(P(X)=\prod_i(X-x_i)\). Then
\[
 P(X-1)=P(X)-\sum_i A_i\prod_{j\ne i}(X-x_j).
\]
The difference of the sides has degree at most \(r-1\) and vanishes at
each of the \(r\) distinct integers \(x_i\), by the defining product for
\(A_i\). A polynomial vanishing at \(u\) is divisible by \(X-u\),
by subtracting its value from each monomial; repeated cancellation
at the other roots proves the required degree bound on the number of
roots. Substitute \(X=1/t\), multiply by \(t^r\), and divide through
by \(\prod_i(1-x_it)\) in the ring \(\mathbb Q[t]/(t^3)\). The factors
are units with inverses \(1+x_it+x_i^2t^2\). This gives the equality
\[
 \prod_i\left(1-\frac{t}{1-x_i t}\right)
   =1-\sum_i A_i \frac{t}{1-x_i t}\pmod{t^3}.
\]
The quadratic coefficient on the left is \(\binom r2-\sum_i x_i\);
on the right it is \(-\sum_i A_i x_i\), as required. For the empty
shape the hook-product expression gives \(W=1\), also with padded
zero rows. The recursion now proves \(W=f^\lambda\) by induction on
\(n\), giving the stated integral identity.

Transposition exchanges the increasing conditions and is its own inverse.
For the rectangle \(x_i=s+r-i\); grouping differences by the smaller
coordinate gives \(\Delta=\prod_{j=0}^{r-1}j!\). For two rows
the shifted coordinates are \(a+1,b\), giving
\((a+b)!(a-b+1)/((a+1)!\,b!)\); expanding the two binomials gives this
same expression.
\end{proof}

\subsection{Semistandard tableaux and the hook-content formula}

\begin{theorem}[Hook-content formula]\label{result16}
For a shape with at most \(r\) rows let \(D_r(\lambda)\) count
semistandard fillings with entries in \(1,\ldots,r\). Then
\[
 D_r(\lambda)=\frac{\Delta(x)}{\prod_{k=0}^{r-1}k!}
 =\prod_{(i,j)\in\lambda}\frac{r+j-i}{h_{ij}}.
\]
For the shape \((5,4,1)\) there are \(288\) standard tableaux and
\(66{,}528\) semistandard tableaux on \(1,\ldots,7\).
\end{theorem}
\begin{proof}
Deleting the entries equal to \(r\) leaves a shape \(\nu\) of at most
\(r-1\) rows. Every such entry is at the right end of its row, in the
sense that all subsequent cells in the row also have entry \(r\).
Every box in row \(r\) must contain \(r\), by column strictness.
The possible shapes are exactly
\(\lambda_i\ge\nu_i\ge\lambda_{i+1}\) for \(1\le i<r\). Indeed the
boxes above row \(i+1\) of \(\lambda\) cannot be deleted. Conversely,
fill the difference of the two shapes with \(r\); the inequalities put
every added box below a retained box if an upper box exists, and place
no box below an added box. This extends every semistandard \(\nu\)
uniquely for that shape. In shifted coordinates
\(y_i=\nu_i+r-1-i\) this says the \(y_i\) range independently over
\(x_{i+1}\le y_i<x_i\). These inequalities already imply strict descent
of the \(y_i\).

For a descending list \(v\) of length \(l\), define
\[
 B_l(v)=\det\bigl(\binom{v_i}{l-j}\bigr)_{1\le i,j\le l}
       =\frac{\prod_{i<j}(v_i-v_j)}{\prod_{k=0}^{l-1}k!}.
\]
For the determinant evaluation, the polynomial \(\binom{X}{k}\) has
leading coefficient \(1/k!\), and the lower-degree monomials are already
spanned by the columns of lower degrees. Subtracting their multiples
successively reduces to the determinant of descending powers with
the displayed factorial factor. The latter determinant is \(\Delta(v)\):
as a polynomial in \(v_1\) it has roots \(v_2,\ldots,v_l\) by equal
rows, and leading coefficient the corresponding smaller determinant
by expansion along row one. The root argument just used above verifies
the factorization for fixed distinct lower entries, and induction gives
the formula. Empty determinants are one.

By induction with \(r-1\) symbols, the deletion count is the sum of
\(B_{r-1}(y)\) over the stated intervals. Expand the determinant as a
signed permutation sum and sum independently in each row. The result
is a determinant of row sums. Pascal's identity telescopes these sums:
\[
 \sum_{z=x_{i+1}}^{x_i-1}\binom zk
   =\binom{x_i}{k+1}-\binom{x_{i+1}}{k+1}.
\]
In the matrix for \(B_r(x)\), subtract row \(i+1\) from row \(i\)
for \(i=1,\ldots,r-1\) in that order. The determinant stays unchanged
by multilinearity and cancellation for repeated rows. Its constant
last column is now zero except for a final one. Expansion there gives
exactly the determinant of the above differences for
\(k=r-2,\ldots,0\). This proves the recurrence, with the empty shape
at \(r=0\) as base. Multiplying the factorials in each numerator row gives
\(\prod_{j=1}^{\lambda_i}(r+j-i)=x_i!/(r-i)!\). The row hook product
proved above gives the content formula. For the example the hook rows
are \(7,5,4,3,1;\ 5,3,2,1;\ 1\), with total product \(12600\).
The numerator is \(10!\) in the standard case and
\((7\cdot8\cdot9\cdot10\cdot11)(6\cdot7\cdot8\cdot9)5\)
in the semistandard case. Division gives the two values.
\end{proof}

\subsection{Permutation insertion and increasing subsequences}

\begin{theorem}[Robinson--Schensted correspondence and Schensted's first-row theorem]\label{result17}
Permutations of \(1,\ldots,n\) correspond bijectively to pairs \((P,Q)\)
of standard tableaux of the same shape \(\lambda\) of size \(n\).
The longest strictly increasing subsequence of the permutation has length
\(\lambda_1\), with length zero for the empty shape. Consequently
\[
  \sum_{\lambda\vdash n}(f^\lambda)^2=n!.
\]
The number of permutations avoiding an increasing subsequence of length
\(k+1\) is the same sum restricted to shapes with \(\lambda_1\le k\).
\end{theorem}
\begin{proof}
Build the insertion tableau \(P\) from the empty filling by reading
the permutation in order. To insert \(a\) in a row, find its leftmost
entry \(b>a\). If found, replace \(b\) by \(a\) and insert \(b\) in the
next row. Otherwise append \(a\) and stop. A missing row acts as empty.
The entries are distinct throughout. A replacement preserves row order,
and no column inequality with a lower neighbor is harmed by decreasing
the cell. If the bump out of a row occurs at column \(c\), its insertion
in the next row occurs at a column \(c'\le c\): the old entry there at
column \(c\), when present, exceeds the bumped value, and a shorter row
can append no later than \(c\). For the new entry at \(c'\) its upper
neighbor is smaller. When \(c'<c\) this follows from the old upper row's
strict order; when \(c'=c\) that neighbor was just replaced by a smaller
value. In particular an appended box in a lower row has an upper box.
The procedure visits at most one row past the old height and adds exactly
one corner. Put the insertion time in this new box in \(Q\). Every
immediately preceding box above or to its left existed earlier, so \(Q\)
is standard.

For the inverse, find the box of the largest remaining time in \(Q\).
It is removable. Delete it in both tableaux, carrying its value \(b\)
from \(P\). In the row above, use the \emph{rightmost} entry \(a<b\),
replace \(a\) by \(b\), and carry \(a\) upward. On passing above row
one, output the carried value. Such a choice always exists when a row
above remains: if the carried value came from column \(c\), that upper
row had a smaller entry at \(c\). The chosen column \(c'\) is therefore
at least \(c\). Replacing \(a\) by \(b\) preserves the row order by
the rightmost rule. Its upper neighbor stays smaller. For its lower
neighbor, the lower row's entry formerly at \(c\) was just deleted or
increased; at a chosen \(c'\ge c\), a lower entry if present now exceeds
the carried \(b\). This verifies all inequalities as the reverse path
proceeds. After any upward replacement, the entries to its left are
below the value \(a\) carried out, so the first entry exceeding \(a\) on
reinsertion is exactly the newly placed \(b\). In the bottom row the
deleted last entry simply appends on reinsertion. In the other direction,
after a forward bump, its replaced entry and all entries to its left
are smaller than the bumped value, and all entries to its right are
larger; the reverse search necessarily picks the same cell. Hence the
local procedures invert each other. Repeating deletions recovers the
permutation in reverse order: each deletion ejects exactly one distinct
label from the remaining set. This proves the bijection and its square sum.

To identify the first row, for a processed prefix call \(m_j\) the
least possible final value of an increasing subsequence of length \(j\),
whenever one exists. The first row after this prefix consists exactly
of these \(m_j\), in increasing order. Prove it by induction, starting
with the empty row. When a new distinct value \(a\) arrives, let \(c\)
be the first column with \(m_c>a\), or the next unused column if there
is none. A subsequence of length \(c\) now ends in \(a\): extend one
ending in \(m_{c-1}<a\) or use the singleton if \(c=1\). No longer
subsequence ending in \(a\) exists, since its prefix of length \(c\)
would have a last element smaller than \(a\), whereas the old \(m_c\)
was absent or larger. For smaller lengths the old minima are already
below \(a\). For greater lengths no new candidate is available. Thus
exactly column \(c\) is replaced or appended, as in the first-row
insertion rule. This proves the claim and the restricted sum.
\end{proof}

\begin{corollary}[Catalan permutation avoidance]\label{result18}
The number of permutations of length \(n\) avoiding an increasing
subsequence of length three is \(C_n\), and the same is true for avoiding
a decreasing subsequence of length three.
\end{corollary}
\begin{proof}
Transpose both tableaux of a pair whose common shape has at most two
columns. A standard tableau with at most two rows gives a length-\(n\)
path with up-steps for top-row labels and down-steps for bottom-row
labels, reading the labels in increasing order. Every prefix has at least
as many up-steps as down-steps, by column strictness. Conversely fill
the next vacant cell of the indicated row at each step; the prefix
condition ensures column strictness. Two such paths have equal final
height exactly when their shapes agree, because their lengths are equal.
Concatenate the first path with the second read backwards and with steps
negated. This gives a Dyck path of length \(2n\), and cutting at its
middle gives the two inverse paths. Reversing the positions in a
permutation interchanges increasing and decreasing subsequences. The
argument also includes \(n=0\).
\end{proof}

\paragraph{Arithmetic bounds and Lean obligations.}
For fixed size bound \(N\) and row count \(r\), take a shape code with
\(r\) digits in base \(\max(2,N+1)\). Tableaux can be stored in
\(rN\) entries, with zero padding; use base \(\max(2,N+1,r+1)\) for
both kinds. Define \texttt{isShape}, \texttt{inShape}, \texttt{isSYT},
\texttt{isSSYT}, \texttt{removable}, \texttt{hook} and
\texttt{shiftedRow} by bounded scans. For the hook identity take
induction on \(m\le N\) with a Boolean fold over all shape codes of
size at most \(m\). The corner proof locates \(m\), erases its cell,
and inverts this transformation by reattaching \(m\).
For \texttt{cornerWeightSum}, implement rational polynomial coefficients
with signed numerator and positive denominator and use arrays through
degree \(r\), plus length-three arrays for the truncated expansion.
The largest shifted coordinate is at most \(N+r\); every factorial in
the formulas is bounded by \((N+r)!\). For the semistandard induction
fix an alphabet bound \(R\) and size bound \(N\) and fold over all shapes
of size at most \(N\) with at most the current number of rows, using
padding to length \(R\). A determinant of size \(l\le R\) is the signed
sum over permutations, each checked in base \(\max(2,l)\). Prove row
linearity by distribution, repeated-row cancellation by swapping the two
row assignments, and cofactor expansion by the chosen column index and
the induced order on the remaining indices. These finite identities
suffice for the telescoping determinant proof.

For insertion take a padded \(N\times N\) tableau. Implement
\texttt{bump}, \texttt{unbump}, \texttt{insert} and \texttt{removeCorner}
with at most \(N+1\) row visits and at most \(N\) comparisons per row.
The local inverse invariant records the carried entry, the previous
column and the untouched prefixes or suffixes of the current and
adjacent rows. Apply the two local inverse lemmas along the stored
row path, then induct on the permutation prefix, retaining the shape
agreement and validity of both tableaux. For the subsequence invariant
enumerate index masks below \(2^N\) to define a least final value or an
absent marker for each length. These are primitive recursive observations
of a fixed finite prefix; the argument uses no quantification over
uncoded maps or infinite series.

\paragraph{Sources.}
The numerical examples appear on the seed's linked
\href{https://en.wikipedia.org/wiki/Young_diagram}{Young diagram/Young tableau}
page. The original papers give the hook formula and the subsequence
interpretation, respectively: J. S. Frame, G. de B. Robinson and R. M.
Thrall, \emph{The Hook Graphs of the Symmetric Group}
(\href{https://doi.org/10.4153/CJM-1954-030-1}{1954}), and
C. Schensted, \emph{Longest Increasing and Decreasing Subsequences}
(\href{https://doi.org/10.4153/CJM-1961-015-3}{1961}).
